\documentclass[11pt]{article}
\usepackage[margin=1in]{geometry}
\usepackage{amsmath,amssymb,amsthm}
\usepackage{booktabs}
\usepackage{rotating}
\usepackage{microtype}
\usepackage[svgnames]{xcolor}
\usepackage[colorlinks=true,linkcolor=Teal,citecolor=Teal,urlcolor=Teal]{hyperref}

\newcommand{\R}{\mathbb{R}}
\newcommand{\C}{\mathbb{C}}
\newcommand{\hd}{d_{\mathrm{head}}}
\newcommand{\nb}{n_b}
\DeclareMathOperator{\softplus}{softplus}
\DeclareMathOperator{\softmax}{softmax}

\title{What the Axes Are Worth:\\
\large measurements on navigation and algorithmic tasks}
\date{}
\author{}

\begin{document}
\maketitle

\begin{abstract}
\noindent
A companion to \emph{Positional Mechanisms in Attention}, which classifies the space
of positional operators and finds two live slots. This paper measures what two
properties of those slots are worth --- properties that no paper in the literature
varies deliberately.

The \textbf{sign} of the phase increment costs nothing and decides what the
accumulator measures: signed increments cancel, so the interval sum is net
displacement and the code is a map; monotone ones cannot, so it is elapsed path and
the code is a clock. On a map task a signed increment beats an index code by
$0.12$--$0.20$ at matched loss on $12/12$ seeds and a monotone one beats it nowhere;
on a task built to want a clock the ordering inverts, and the same unconstrained
architecture learns a different accumulator for each. The \textbf{input-side rank}
of the content-to-phase map is a step rather than a slope: $r=2$ learns a badly
conditioned basis --- opposite actions fail to cancel by half the action scale --- and
$r=4$ repairs it for $384$ parameters.

Both are measured on grid navigation and algorithmic tasks rather than on text,
because they are separable only where position must be reconstructed from the input
rather than read off the index. Every evaluation is on an environment redrawn at test
time, which makes them transfer measurements: a structural code reaches $0.99$ where
a parameter-matched index code sits on the $0.506$ blank floor.

We also report what did not work. An explicit content gate separates action from
observation tokens by a factor of four and improves accuracy nowhere; a published
external benchmark does not discriminate the mechanism, for a reason our own
per-offset data predicted; and the growth exponent that unifies the sign and rank
results is nearly collinear with a simpler observable and does not, on the evidence
here, explain either.
\end{abstract}

\section{What is measured, and how}
\label{sec:setup-b}

The axes referred to throughout --- A1 factor, A2 locus, A3 driver, A4 algebra, A5
sign, A6 rank, A7 host --- are those of the companion review, and the relational map
placing every mechanism on them is Table~1 there. This paper assumes that
classification and tests two of its cells.

Standard errors are over seeds; \textsc{mde} is $2.8\,\mathrm{sd}/\sqrt{n}$, and a
contrast inside its \textsc{mde} is reported as \emph{unmeasured} rather than as a
null. Contrasts are paired per seed, and arms compared against one another were
trained in one batch. Where accuracy correlates with final training loss --- it has
been as strong as $r=-0.999$ --- contrasts are given raw and loss-matched, and both
are reported.

\section{Factorisation, and what the separation buys}
\label{sec:factorisation}

TEM's thesis is a factorisation: structural knowledge $g$ is environment-invariant,
sensory input $x$ is environment-specific, and memory stores the conjunction
$p = g \otimes x$. Transfer follows because $g$ survives a change of environment
while $x$ does not.

The results above say that MapFormer is that factorisation carried into attention
with a parallel scan. The accumulated phase $\theta_t = \omega\sum\Delta_u$ is $g$:
identical in any environment with the same transition structure, because it depends
on the action stream and not on what was seen.

\textbf{But that separation is learned, not built} --- the difference between a
factorisation the architecture guarantees and one it merely tends to find.
MapFormer's increment is computed from \emph{every}
token embedding --- $\Delta_u = W_{\mathrm{out}}W_{\mathrm{in}}x_u$ with no gating by
token type --- so the model has to \emph{discover} that $\Delta \approx 0$ on
observations. It does: on a trained model actions move the phase about five times as
much as observations do. And the recency result of $\S$\ref{sec:crossover} is the
same phenomenon isolated --- an unconstrained increment learns a content gate in
$8/8$ seeds, and a magnitude-matched intervention shows the gate is most of the
mechanism rather than a component of it. So the separation is real and measurable;
it is emergent rather than enforced.

\textbf{And that is the design, not an oversight.} MapFormer's own text is explicit:
``to learn a cognitive map, the model has to identify that action tokens update the
internal position, while ensuring that observations leave the structure untouched.
This is achieved through a two-stage projection process'' --- the low-rank
$W_{\mathrm{out}}W_{\mathrm{in}}$ \emph{is} the identification mechanism, with
$W_{\mathrm{in}}$ mapping each token to ``a low-dimensional action representation
$\Delta_t \in \R^r$ \dots capturing the actual motion to perform''. Hard-coding the
split would remove exactly what the mechanism is for, and would confine it to
settings with a labelled action stream --- which excludes text, where there are no
action tokens and RoPE is the $\Delta\equiv1$ special case.

This has a consequence for $\S$\ref{sec:rankmeas} worth stating: \textbf{the
bottleneck is the separator}. The rank $r$ is not merely the dimension of a
displacement code, it is the channel through which the what/where distinction has to
pass, so the finding that $r=2$ learns a skewed basis and $r=4$ repairs it for $384$
parameters is a finding about the \emph{separation} and not only about the geometry.
An enforced split is therefore an oracle \emph{control} rather than a proposal: it
upper-bounds what the learned separation can reach, and the gap measures what
learning the split costs. If the two match, the design is vindicated with a number. Content is $x$. And the conjunction
is exactly the WM/EM distinction --- additively in the score for MapFormer-WM, and
for MapEM \emph{literally} the tensor product, since
$(A_X \odot A_P)_{ts} = (q^x_t \otimes q^p_t)\cdot(k^x_s \otimes k^p_s)$. TEM's
$g \otimes x$ and MapEM's Hadamard product are the same object; what MapFormer adds
is that the interval operator composes by prefix sum, so the where can be built in
$O(\log T)$ instead of recurrently.

\paragraph{Why the measurements are transfer measurements.} This is easy to miss in
the tables above and is the reason they are set up as they are. \textbf{Every
evaluation here is on an environment the model has never seen}: the observation map
is redrawn from a held-out seed, so the entire what-to-where binding is new, and
only structure can carry over. Against a measured always-predict-blank floor of
$0.506$:

\begin{table}[htbp]\centering\small
\caption{Transfer to a new environment with the same relational structure. In every row the sensory map is redrawn at evaluation, so nothing environment-specific survives; the gap is what the structural code carries.}
\label{tab:transfer}
\begin{tabular}{@{}p{4.6cm}p{3.4cm}p{3.0cm}p{3.2cm}@{}}
\toprule
task & structural code & index code & floor \\
\midrule
torus revisit, fresh map ($n{=}8$) & $\mathbf{0.994}$ / $0.987$ / $0.967$ & $0.530$ / $0.509$ & $0.506$ \\[2pt]
blind continuation, $64^2$ ($n{=}5$) & $\mathbf{0.730 \pm 0.247}$ & $0.154 \pm 0.018$ & $0.063$ \\[2pt]
blind continuation, $128^2$ ($n{=}3$) & $\mathbf{0.823 \pm 0.043}$ & $0.192 \pm 0.022$ & $0.063$ \\[2pt]
compositional motif transfer ($n{=}8$) & $\mathbf{0.415 \pm 0.096}$ & $0.216 \pm 0.004$ & $\approx 0.072$ \\
\bottomrule
\end{tabular}
\end{table}

\noindent An index code sits at or near the floor on a fresh map. It is not that it
predicts badly --- it cannot locate itself, so it has nothing to bind an observation
\emph{to}, and the map it built in the previous environment is worthless. A
structural code arrives with the where already valid and needs only to rebind. That
is the factorisation claim, measured.

\paragraph{The separation is what is doing the work, not the architecture.} Two
controls make this attributable. The path-integrated and index arms here are
parameter-matched to within $0.4\%$ and differ only in whether the phase is driven
by the action stream or the index, so the gap is the \emph{driver}, not capacity.
And the blind-continuation result survives context destruction --- shuffling the
explore-phase observations takes it from $0.918$ to $0.074$ --- so the model is
using the conjunction it built in context and not a statistical regularity of the
task.

\paragraph{Two boundaries, both of which the factorisation predicts.} The separation
buys nothing when the structural code cannot be \emph{formed}: under rotation-based
actions the displacement depends on accumulated heading, a fixed per-token increment
cannot represent it, and the measured effect falls from $+0.438$ to $+0.050$ --- an
index code then does as well, because neither has a usable where. And it buys
nothing when there is no structure to share: the effect is a threshold in map extent
($-0.010$, $+0.015$, $+0.305$ for $32$, $128$, $512$ occupied cells), so on a small
environment, where an episode revisits everything anyway, there is nothing for a
structural code to generalise \emph{from}.

\paragraph{What is not shown.} These are transfers to a new \emph{instance} of the
same structure --- a redrawn observation map on the same torus, new motif instances
in the same grammar. Transfer across a change of \emph{structure} --- a different
topology, a different action space, or a graph learned in one environment reused in
another --- is the claim TEM makes most strongly and is not tested here. So is
\emph{faster} learning: everything above is final accuracy, and no sample-efficiency
curve has been measured. Both are the obvious next experiments, and
$\S$\ref{sec:predictions} lists what the frame predicts for them.


\section{Measurements}
\label{sec:measured}

The axes above are separable only where position must be reconstructed from the input
rather than read off the index, which is why the measurements below are on
navigation and algorithmic tasks rather than on text.

\paragraph{Setup.} The primary task is \textbf{torus revisit prediction}: an agent
takes actions on an $N\times N$ torus, observations are aliased across cells ($16$
types over $N^2$ cells, half of them blank), the token stream interleaves actions and
observations, and the model predicts the observation at \emph{revisited} cells only.
Unless stated otherwise: $N=64$, one layer, two heads, $d_{\mathrm{model}}=128$,
$\nb=32$, trained at sequence length $128$ for $300$ epochs with warmup and cosine
decay at $\mathrm{lr}=10^{-3}$, and evaluated on a held-out observation map.
``Accuracy'' is next-token accuracy at revisited positions, against a measured
always-predict-blank floor of $0.506$. Two further tasks appear where noted:
\textbf{parity} at sequence length $L$, and \textbf{blind continuation}, in which the
model explores with observations revealed and then predicts the observation at each
cell with observations withheld (chance $1/16 = 0.0625$).

Standard errors are over seeds; \textsc{mde} is $2.8\,\mathrm{sd}/\sqrt{n}$, and a
contrast inside its \textsc{mde} is reported as \emph{unmeasured} rather than as a
null. Contrasts are paired per seed, and arms compared against one another were
trained in one batch.

\paragraph{What these measurements assert.} Three findings and one negative, each
with the contrast that would have refuted it.

\begin{itemize}\itemsep2pt
\item \textbf{F1. The sign of the increment is load-bearing where the task needs a
  map, and free.} At matched loss and identical parameter count, a signed increment
  beats an index code by $0.12$--$0.20$ ($12/12$ seeds); a monotone one beats it
  nowhere ($\S$\ref{sec:sign}). \emph{Refuted by} a monotone arm matching the
  signed arm on the map task.
\item \textbf{F2. Rank $r=2$ fails by conditioning, not capacity.} $r=4$ buys
  $+0.085$ for $384$ parameters and is flat to $r=32$; the learned code at $r=2$ is
  skewed rather than small ($\S$\ref{sec:rankmeas}). \emph{Refuted by} a capacity
  slope, or by the deficit growing with world dimension --- which was tested, and
  it does not.
\item \textbf{F3. One quantity accounts for F1 and F2, and it reads off the
  task.} The accumulator's growth exponent $\alpha$ correlates with cancellation at
  $r = +0.9995$ across arms ($\S$\ref{sec:growth}), and an unconstrained
  architecture moves it from $0.591$ to $0.967$ between a map task and a clock task
  while every constrained arm stays fixed ($\S$\ref{sec:crossover}).
  \emph{Refuted by} the free arm learning the same accumulator on both.
\item \textbf{N1. The axis on which most mechanisms help is not explained.} Four
  separate mechanisms help at extrapolated length; $\alpha$ accounts for two and
  demonstrably not the others, and the imported long-context account is refuted
  outright ($\S$\ref{sec:localisation}). This is the gap, stated as one.
\end{itemize}

\noindent The remaining subsections are supporting rather than load-bearing: what
the neighbouring constraints are worth ($\S$\ref{sec:knobs}), why these axes are
invisible on text ($\S$\ref{sec:navigation}), and a cost/benefit summary
($\S$\ref{sec:anchors}).

\subsection{F1: the sign is load-bearing, and free}
\label{sec:sign}

Six arms $\times$ $12$ seeds, one batch, $r=4$ throughout, \textbf{parameter count
identical at $204{,}757$} --- the constraint adds nothing and does not touch the rank.

\begin{table}[htbp]\centering\small
\caption{The sign ablation, \textbf{raw} accuracy. Six arms $\times$ $12$ seeds, one batch, identical parameter count ($204{,}757$). The construction control \texttt{Vanilla\_r4}, omitted here for width, is within $+0.006$ of \texttt{signed} at every length --- \emph{unmeasured}, which is the condition for the batch to be readable at all. Read this table beside the loss-matched contrasts below: raw and loss-matched point the same way for the signed-vs-monotone comparison and opposite ways for monotone-vs-index.}
\label{tab:sign}
\begin{tabular}{llrrrr}
\toprule
arm & increment & $T{=}128$ & $T{=}512$ & $T{=}1024$ & final loss \\
\midrule
signed  & $W_{\mathrm{out}}W_{\mathrm{in}}x$ & $1.000$ & $0.978$ & $0.922$ & $0.0002$ \\
\midrule
$|\cdot|$ & $|W_{\mathrm{out}}W_{\mathrm{in}}x|$ & $0.946$ & $0.675$ & $0.558$ & $0.171$ \\
softplus  & GRAPE-AP form                       & $0.977$ & $0.798$ & $0.584$ & $0.068$ \\
CARoPE    & $1/(\softplus(\cdot)+1)$            & $0.900$ & $0.809$ & $0.645$ & $0.293$ \\
\midrule
RoPE      & index                               & $0.799$ & $0.449$ & $0.345$ & $0.784$ \\
\bottomrule
\end{tabular}
\end{table}

\noindent The raw column is not the contrast of interest, and says so plainly: the
monotone arms fit far worse ($0.171$, $0.068$, $0.293$ final loss against $0.0002$),
so a raw comparison against an index arm that fits worse still ($0.784$) reports
convergence rather than representation. \textbf{At matched training loss, signed path
integration beats an index code by $+0.123$ and $+0.195$ at $T=512$ and $T=1024$
($12/12$ seeds), while monotone path integration beats it nowhere} --- $-0.028$ at
$T=128$ and unmeasured beyond. Note also that the index arm sits at or below the
measured always-predict-blank floor of $0.506$ at both OOD lengths, so its raw
margins are not interpretable as skill. The entire measured value of a content-dependent phase over a plain index
clock, on this task, requires the increment to be signed.

The learned code shows why. The opposition score
$\lVert\Delta(+x)+\Delta(-x)\rVert / \overline{\lVert\Delta\rVert}$ is $0$ when
opposite actions cancel exactly and $2$ when they are identical:

\begin{table}[htbp]\centering\small
\caption{The learned action code. Opposition is $0$ when opposite actions cancel exactly and $2$ when they are identical.}
\label{tab:opposition}
\begin{tabular}{lrrr}
\toprule
arm & oppose$_x$ & oppose$_y$ & $|\cos(N,E)|$ \\
\midrule
signed    & $\mathbf{0.125}$ & $\mathbf{0.106}$ & $0.218$ \\
$|\cdot|$ & $1.849$ & $1.855$ & $0.587$ \\
softplus  & $1.905$ & $1.885$ & $0.723$ \\
CARoPE    & $\mathbf{1.981}$ & $\mathbf{1.977}$ & $\mathbf{0.930}$ \\
\bottomrule
\end{tabular}
\end{table}

\noindent CARoPE's published parameterisation reaches $1.98$ of a possible $2.00$,
with its north and east axes all but collapsed onto one another. These models do not
merely score worse --- they cannot represent a $-1$ action, and the probe confirms
they do not.

At the training length the signed baseline is at $1.000\pm0.000$,
so accuracy there cannot show a deficit of any size. The training-length effect is
visible in the \emph{loss} instead: every constrained arm fits strictly worse on
$12/12$ seeds at identical parameters.

This is a replication in a new regime. The constraint and its
consequence for parity are established (\Sthe companion review(ii)); what is new is
navigation rather than formal languages, and an isolation in which $|\Delta|$ against
$\Delta$ is a single operation at identical parameter count.

\subsection{F2: rank $r=2$ fails by conditioning, not capacity}
\label{sec:rankmeas}

The bottleneck $r$ is meant to be the dimensionality of the action space --- for a
$2$D grid, a $2$-vector --- which makes $r=2$ sufficient \emph{by construction}. It is
not sufficient in practice. Against $r=2$ at $8\times$ the training length ($T=1024$
against a training length of $128$; $8$ seeds):

\begin{table}[htbp]\centering\small
\caption{The rank sweep against $r=2$, $8$ seeds. A step at $r=2$, flat from $r=4$ to $r=32$.}
\label{tab:rank}
\begin{tabular}{lrrr}
\toprule
 & params added & $T{=}512$ & $T{=}1024$ \\
\midrule
$r=4$  & $+384$    & $+0.038$ & $\mathbf{+0.085}$ ($8/8$, $t\,3.57$) \\
$r=8$  & $+1{,}152$ & $+0.037$ & $+0.091$ \\
$r=16$ & $+2{,}688$ & $+0.028$ & $+0.079$ \\
$r=32$ & $+5{,}760$ & $+0.038$ & $+0.095$ \\
\bottomrule
\end{tabular}
\end{table}

\noindent This is a \textbf{step at $r=2$, not a capacity slope}: flat from $r=4$ to
$r=32$. The cause is conditioning, not capacity. Given four dimensions the model still
places its actions in a $2$-plane ($100.0\%$ of spectral energy; $99.96\%$ even at
$r=32$), vindicating the dimensional argument --- what fails is the \emph{basis} that
$r=2$ is forced into:

\begin{table}[htbp]\centering\small
\caption{The learned basis at $r=2$ and $r=4$. The deficit is conditioning, not capacity.}
\label{tab:basis}
\begin{tabular}{lrrr}
\toprule
 & opposition ($N{+}S\approx0$) & $|\cos(N,E)|$ & obs/action norm \\
\midrule
$r=2$ & $0.495$ & $0.783$ & $0.139$ \\
$r=4$ & $\mathbf{0.092}$ & $\mathbf{0.175}$ & $0.042$ \\
\bottomrule
\end{tabular}
\end{table}

\noindent At $r=2$ opposite actions fail to cancel by half the action scale and north
and east are nearly parallel. The seed standard deviation also falls from $0.064$ to
$0.012$ at $T=1024$, which matters because it sets every detectable effect size.

Two further observations. The deficit does \emph{not} grow with the dimensionality of
the world: at $D=5$, $r=2$ is unimpaired ($0.896$, its best of three conditions), so
the packing account of the companion review does not explain it. And an unconstrained
generator cannot be projected below rank $\approx16$ post hoc while $r=4$ trains fine
--- trained-at-rank and truncated-to-rank are different objects.

\subsection{Supporting: what the neighbouring constraints are worth}
\label{sec:knobs}

Selective RoPE's three additions to MapFormer's generator, measured on the torus at
$T=1024$ and on a parity task at $L=128$; increments against MapFormer at $r=2$.

\begin{table}[htbp]\centering\small
\caption{Selective RoPE's additions to MapFormer's generator, against MapFormer at $r=2$.}
\label{tab:knobs}
\begin{tabular}{lrll}
\toprule
constraint removed & params & parity $L{=}128$ & torus $T{=}1024$ \\
\midrule
$K=1 \to 4$ (causal conv)      & $+193$     & $-0.020$ & $-0.064$ ($t\,{-}1.99$) \\
$m = r \to H\nb$ (full rank)   & $+7{,}873$ & $-0.020$ & $+0.058$ ($7/8$) \\
$\varphi$ linear $\to$ gated   & $+8{,}193$ & $-0.030$ & $+0.086$ ($8/8$) \\
all three                      & $+16{,}385$& $-0.009$ & $+0.048$ ($t\,1.36$) \\
\midrule
\emph{$r=2\to4$, constraint kept} & $\mathbf{+384}$ & --- & $\mathbf{+0.085}$ ($8/8$) \\
\bottomrule
\end{tabular}
\end{table}

\noindent Two readings. First, \textbf{the two $\approx$8k-parameter knobs are
statistically indistinguishable from each other} ($+0.018$ and $+0.028$, both inside
their \textsc{mde}s), so what they buy is most plausibly capacity rather than either
named mechanism --- and widening the bottleneck MapFormer already has buys the same
thing for $384$ parameters, $21\times$ cheaper. Second, the convolution is the only
knob negative on both tasks, consistent with the companion review: $K>1$ costs the
homomorphism.

These arms replace the readout $\operatorname{diag}(\omega)W_{\mathrm{out}}$
by $\tau W_2$. For the conv and gate rows the function class is preserved exactly ---
it is the construction of the companion review --- so the difference is initialisation
(a geometric frequency ladder against a default draw) plus one scalar; those effects
are attributable up to initialisation. The two full-rank rows also add weight
normalisation and a bias, so their function class genuinely differs.

A separate probe rules out the natural mechanism for the gate. If the sigmoid worked
by suppressing observation tokens --- which MapFormer must instead learn to send to
$\Delta\approx0$ --- the gate would be near zero on observations. It is $0.416$, and
the action/observation ratio is $1.35\times$ on the torus where the gate helps and
$1.54\times$ on parity where it hurts.

\subsection{Supporting: why these axes are invisible on text}
\label{sec:navigation}

Every paper in the companion review evaluates on language, recall, state tracking, or ---
for LieRE and Algebraic PE --- images, video and trees. None runs grid navigation, and
it is the regime in which these axes are distinguishable at all.

The basic separation is large. On a task where the model explores with observations
revealed and then continues blind, predicting the observation at each cell, path
integration scores $0.730 \pm 0.247$ ($n=5$, $64^2$) against $0.154 \pm 0.018$ for an
index code, chance $0.0625$, with no seed overlap. The axis is path integration and not
the encoding: PoPE with path integration scores $0.847$ against $0.117$ for PoPE with
an index (both $n=3$; only the MapWM pair was extended to $n=5$, and it fell $0.158$
when it was, so read the ratio and not the values). On the standard revisit task, re-derived from the per-seed data
($n=8$, held-out map), the position main effect is $+0.461$ (8/8 seeds, MDE
$0.027$) and the encoding main effect is $+0.003$, which is below its own MDE of
$0.029$ at 5/8 seeds. No ratio is quoted: with a denominator that close to zero the
quotient is not a stable quantity. The defensible statement is that one factor
moves the result by roughly half and the other does not move it measurably.

Two boundaries are worth stating. The effect is not a property of navigation as such
but of \emph{map extent}: at matched observation aliasing it is $-0.010$, $+0.015$ and
$+0.305$ for $32$, $128$ and $512$ occupied cells --- a threshold, not a gradient
($n=3$ per condition; the $512$ endpoint that makes it a threshold is the least
replicated point, though the $0.290$ jump is roughly twice the measured noise floor). And
it does not survive rotation-based action spaces: under turn/turn/forward, where
displacement depends on accumulated heading, the effect falls from $+0.438$ to $+0.050$,
which is the companion review's commutativity limitation made concrete. Recording the
absolute displacement instead of the commanded turn restores it to $+0.488$.

\subsection{Supporting: what each ingredient costs and buys}
\label{sec:anchors}

\begin{table}[htbp]\centering\small
\caption{Anchors: what each ingredient costs and buys on the torus task.}
\label{tab:anchors}
\begin{tabular}{p{2.4cm}p{1.9cm}p{2.6cm}p{6.4cm}}
\toprule
mechanism & sets & parameter cost & measured effect \\
\midrule
Path integration & phase $\theta$ & $+448$ ($+0.23\%$) & $+0.078$ parity ($16/16$); $+0.461$ torus \\[2pt]
\textbf{Signed} increment & sign of $\Delta$ & $\mathbf{0}$ & largest per parameter; removing it costs $-0.215$/$-0.280$ at $T{=}512/1024$ \\[2pt]
Bottleneck rank & state rank $\rho$ & $+384$ ($+0.19\%$) & $+0.085$ at $T{=}1024$ for $r{=}2\to4$ ($8/8$); flat to $r{=}32$ \\[2pt]
PoPE & magnitude $\mu$ & $+320$ & $+0.027$ \emph{with} path integration; both arms on the floor without it \\[2pt]
\bottomrule
\end{tabular}
\end{table}
\subsection{N1: the length axis, and a refuted import}
\label{sec:localisation}

Transposed to a content-dependent
phase the coordinate is the accumulator rather than the index, which gives a
prediction: the benefit of the OOD-helping mechanisms should localise to the
channels whose period exceeds the trained range of $S_t$, and ablating those
channels should cost \emph{less} at extrapolated length because they are already
being read out of distribution.

The precondition holds --- at the training length $8$--$17$ of $64$ channels never
complete a cycle of their own period, and at $8\times$ that length only $1$--$4$ do
--- so the mechanism could operate. It does not. Ablating the $32$ lowest-frequency
channels is free at training length ($-0.001$) and costs $0.14$--$0.25$ at $T=1024$,
on every arm whose baseline is off the floor. \textbf{Those channels are more
load-bearing at extrapolated length, not less}, which is the opposite of the
prediction. Damage is localised in frequency --- high-frequency ablation costs
$0.42$--$0.57$ --- but the plainer reading wins: the low-frequency channels carry
position at map scale, and coarse position matters more when the walk covers more of
the map.

\subsection{F3: one quantity accounts for both}
\label{sec:growth} All
four arms start from the same accumulator range at training length ($88$--$94$) and
separate entirely by how fast it grows with sequence length,
$\operatorname{range}(S) \sim T^{\alpha}$:

\begin{table}[htbp]\centering\small
\caption{The accumulator's growth law. All four arms start from the same range at training length and separate by exponent.}
\label{tab:growth}
\begin{tabular}{lrrrr}
\toprule
arm & $T{=}128$ & $T{=}1024$ & $\alpha$ & opposition score \\
\midrule
signed, $r{=}4$   & $91.8$ & $269.7$ & $\mathbf{0.518}$ & $0.125$ \\
signed, $r{=}4$ (rank sweep) & $94.1$ & $279.7$ & $\mathbf{0.524}$ & $0.092$ \\
signed, $r{=}2$   & $88.5$ & $320.5$ & $0.619$ & $0.495$ \\
monotone, $r{=}4$ & $91.5$ & $649.8$ & $\mathbf{0.943}$ & $1.849$ \\
\bottomrule
\end{tabular}
\end{table}

\noindent A signed, well-conditioned increment gives $\alpha \approx 0.52$: the
accumulator performs a \textbf{diffusive random walk}. A monotone one gives
$\alpha \approx 0.94$: \textbf{ballistic drift}. The last column is the opposition score of
$\S$\ref{sec:sign}; across the four arms $r(\text{opposition},\alpha) = +0.9995$.

\paragraph{What $\alpha$ is worth.} It is an observable rather than a knob, and it
is nearly collinear with the opposition score ($r = +0.9995$), which is read
directly off the learned action code. So it carries little information opposition
does not, and its contribution is economy: it makes the sign and rank results one
finding at two severities rather than two. One use is its own --- opposition needs
labelled opposite actions and cannot be computed on text at all, while $\alpha$
needs only trajectories, so it is the form this measurement takes in the setting
this literature works in. That measurement has not been made.

\textbf{This is one mechanism for two separate results.} How completely opposite
actions cancel sets whether the accumulator diffuses or drifts; that sets how fast it
leaves the range it was trained on; and that co-varies with the rate of degradation
with length. The last link is the weak one and is stated as association rather than
cause: $\theta$ enters through $\cos,\sin$, so a larger accumulator is not
\emph{prima facie} out of distribution at all, and the standard account of why it
would be --- untrained low-frequency channels --- is refuted here
($\S$\ref{sec:localisation}). A skewed rank-$2$ basis (opposition $0.495$) and a monotone increment
($1.849$) are the same failure at different severities. The relation is correlational across four
arms, and $\alpha\approx0.5$ for a signed walk is unsurprising in itself; what it
rests on is the spread to $0.94$ and the agreement with an independently measured
quantity.


\subsection{F3, continued: the crossover}
\label{sec:crossover}

Every measurement above is on a task that wants a \emph{map}. That makes
``a signed increment is better'' and ``a signed increment suits this task''
indistinguishable, and it is the obvious objection to the companion review. The
test is a task whose answer depends on elapsed count rather than net displacement.

\textbf{The task.} A stream of symbols carries interleaved queries; a query names
an offset $k$ and the model must retrieve the $k$-th most recent \emph{content}
symbol, where uncounted filler tokens are interleaved at rate $0.5$. It is a
retrieval, not a decode --- asking a MapFormer to \emph{report} elapsed time
reproduces a known failure, since the code makes positions comparable and not
readable. The filler is what makes $k$ a \emph{contextual} position in CoPE's
sense: the token distance back to the answer is $129.7 \pm 10.3$ at $k=64$, so no
fixed offset addresses it. Without the filler the task is index retrieval in
disguise, and measured that way an index code scores $1.000$ while signed and
monotone tie. Six arms $\times$ eight seeds, one batch, $k_{\max}=64$, trained at
$T=1024$; chance $0.0625$, shortcut gates at chance, and doubling the budget moves
the weakest arm by $-0.009$.

\begin{table}[htbp]\centering\small
\caption{The crossover. The cost of constraining the increment to be monotone, paired per seed, on a task that wants a map and one that wants a clock.}
\label{tab:crossover}
\begin{tabular}{lrrl}
\toprule
 & torus (a map) & recency (a clock) & \\
\midrule
monotone $-$ signed, raw & $-0.280$ & $-0.004$ & \\
monotone $-$ signed, loss-matched & $-0.215$ & $+0.026$ & both inside MDE \\
seeds & 12/12 negative & 4/8 & \\
\bottomrule
\end{tabular}
\end{table}

\noindent Giving up cancellation costs a fifth to a quarter of accuracy on the map
task and \textbf{nothing measurable} on the clock task: an interaction of about
$+0.28$. The accumulator says the same thing more sharply. The unconstrained arm
may adopt either code, and adopts a different one per task, while every arm that
\emph{cannot} choose sits at the same exponent on both:

\begin{table}[htbp]\centering\small
\caption{The growth exponent is a readout of what the task demanded. Only the unconstrained arm can move, and it does.}
\label{tab:alphashift}
\begin{tabular}{lrrr}
\toprule
arm & $\alpha$ (torus) & $\alpha$ (recency) & $\Delta$ \\
\midrule
\textbf{signed, unconstrained} & $\mathbf{0.591 \pm 0.028}$ & $\mathbf{0.967 \pm 0.009}$ & $\mathbf{+0.376}$ \\
$|\Delta|$ (monotone) & $1.010$ & $0.976$ & $-0.033$ \\
$\mathrm{softplus}$ (monotone) & $1.005$ & $1.005$ & $+0.000$ \\
CARoPE form (monotone) & $1.003$ & $0.992$ & $-0.011$ \\
\bottomrule
\end{tabular}
\end{table}

\noindent $n=12$ and $8$; the shift has a standard error of $0.009$. The
constrained rows are the control that attributes it to the \emph{choice} rather
than to the two tasks' statistics.

This is corroboration rather than the finding: the result is the accuracy
interaction above, and the $\alpha$ table measures the same thing on the learned
code instead of on held-out accuracy. Worth having because the two could have
disagreed; not independent evidence.

\paragraph{What supplies the clock, established by intervention.} The arm does not
make every increment positive --- the negative fraction of $\Delta$ barely moves,
$0.498 \to 0.478$. It learns a \textbf{content gate}: large increments on counted
tokens, near-zero on filler, present in $8/8$ seeds per head. Gate \emph{strength}
does not predict the outcome ($r=-0.34$ across a near-ceiling range), so this was
settled by eval-only intervention rather than by more measuring. Removing the
filler increments is a no-op ($1.000 \to 0.9997$); removing the counted ones
destroys the task ($0.068$), which is the positive control. The decisive contrast
holds magnitude fixed: replacing $\Delta$ with a \emph{constant} on content alone
scores $0.783$, and with a constant on \emph{every} token --- same total drift,
differing only in which tokens it lands on --- $0.189$, a gap of $+0.594$ at $8/8$
seeds. A constant increment on content recovers most of the baseline, so the gate
is most of the mechanism rather than a component of it.\footnote{An earlier
condition that simply made the filler count also collapsed, apparently decisively.
It is confounded and is not counted: a control that changed \emph{only} the
accumulator's scale, still counting content alone, collapsed just as hard. This
model is acutely sensitive to $\theta$'s magnitude, so any intervention that moves
it is destructive for reasons unrelated to counting. The magnitude-matched pair was
added afterwards, as a control for that confound.}

\paragraph{And a fixed index code cannot do it at all.} The path-integrated arms
reach $0.96$--$1.00$ and the index arms $0.234$, a gap of $+0.750$ at $8/8$ seeds
--- larger than the navigation effect of $\S$\ref{sec:navigation}. The per-offset
curve says why: the index arms solve $k=1$ ($0.99$) and collapse from $k=8$ on
($0.11$--$0.33$), while every path-integrated arm is flat out to $k=64$. This
reproduces CoPE's own claim in a different architecture --- against contextual
position, a relative code's best is a decaying attention --- and is cited here as
an anchor rather than as a new result. Note also that a token clock built inside
this architecture ($0.189$, above) lands \emph{below} the index arms, which is
what it is.

\subsection{Where MapPoPE lands}
\label{sec:mappope}

This review began from a concrete question --- MapFormer supplies a
content-dependent \emph{phase}, PoPE a content-dependent \emph{magnitude}, and the
coordinates are disjoint, so what does imposing both buy? --- and the frame exists
to answer it in a way that generalises. The answer is mixed, and the shape of the
mixture is the useful part.

\begin{table}[htbp]\centering\small
\caption{MapPoPE's record. It is the strongest arm on the task the phase mechanism was designed for, and loses to a pure index code on the task where path integration itself fails.}
\label{tab:mappope}
\begin{tabular}{@{}p{4.5cm}p{3.1cm}p{3.3cm}p{3.4cm}@{}}
\toprule
task & MapPoPE & best alternative & reading \\
\midrule
torus revisit, held-out map ($n{=}8$) & $\mathbf{0.994 \pm 0.017}$ & MapEM-os $0.987$; MapFormer $0.967$ & \textbf{best measured} \\[2pt]
torus, OOD $l = 2048$ ($n{=}8$) & $\mathbf{0.970 \pm 0.028}$ & MapFormer $0.854$ & \textbf{best measured} \\[2pt]
PoPE's contribution, loss-matched, 4 grids ($n{=}8$) & $+0.077$ to $+0.095$ & --- & helps; \emph{flat} in octaves, so the frequency-count account is refuted \\[2pt]
MiniGrid DoorKey-16$\times$16, $T{=}1024$ ($n{=}8$) & $0.942 \pm 0.017$ & \textbf{PoPE + \emph{index}} $0.955 \pm 0.003$ & \textbf{loses to an index code} \\[2pt]
compositional transfer ($n{=}8$) & $0.429 \pm 0.117$ & MapFormer+hier $0.415 \pm 0.096$ & tied \\[2pt]
enwik8 bpc ($n{=}3$) & $1.3786$ (lowest) & PoPE+index $1.3806$ & numerically best, \textbf{not established}: the margin is under its own checkpoint sd \\
\bottomrule
\end{tabular}
\end{table}

\noindent Three things follow, and the third is the one worth carrying away.

\textbf{It is the best arm on the task the phase mechanism was built for}, at the
training length and at $16\times$ it. That is a real result and it is where the
combination earns its place.

\textbf{It is not a general improvement.} On MiniGrid it is beaten by PoPE with an
\emph{index} phase, because rotation-based actions defeat a fixed per-token
increment ($\S$\ref{sec:navigation}) --- so the phase half is a liability exactly
where path integration itself is. And on text the margin is inside its own noise.
A combination inherits its components' failure modes, not only their strengths.

\textbf{It fills a cell in the table without supplying what that cell promises.}
MapPoPE reads as ``both slots'' in the companion review, and the companion review
says a system needing map \emph{and} clock must spend both. But PoPE's magnitude is
\emph{instantaneous} --- read off the current token, not accumulated --- so it adds
no second accumulator, and measured, it does not move the growth exponent at all
($\alpha = 0.689$ against a baseline $0.710$, inside its \textsc{mde}). MapPoPE is
two mechanisms in the phase-and-magnitude slots that supplies \emph{one} answer.
The genuinely unoccupied configuration is a signed accumulated phase beside an
\emph{accumulated} content-dependent magnitude in softmax attention, and
$\S$\ref{sec:predictions} says what to build there.


\section{Predictions}
\label{sec:predictions}

The frame is worth something only if it says what to try next. These follow from
the axes rather than from the measurements, and none has been run. Each is stated
with what would refute it, so that the section is falsifiable rather than a wish
list.

\subsection{Mechanisms on tasks they have not been evaluated on}

\begin{table}[htbp]\centering\small
\caption{Untested deployments the frame implies. None of these experiments exists.}
\label{tab:predictions}
\begin{tabular}{@{}p{2.9cm}p{3.2cm}p{4.2cm}p{3.5cm}@{}}
\toprule
mechanism & task & why the frame predicts it & refuted by \\
\midrule
PaTH, DeltaNet, RWKV-7 & navigation with \textbf{rotation} actions & non-abelian transport is exactly what turn/forward needs, and a fixed per-token increment cannot represent it: MapFormer falls from $+0.438$ to $+0.050$ there & landing at MapFormer's rotate level, which would mean the deficit is not about commutativity \\[3pt]
MapFormer & \textbf{state tracking} ($A_5$, $S_5$), formal languages & a signed increment is a parity register --- measured $+0.316$ at $L{=}16$ --- and it keeps the parallel scan, which the delta-rule family forfeits & failing where DeltaNet succeeds, at matched budget \\[3pt]
Mamba-3 & \textbf{cognitive-map} navigation & the only mechanism holding both slots with a signed rotation; should be the strongest non-attention arm & falling to Mamba-2/GLA's level, which would mean the rotation is not what navigation uses \\[3pt]
LieRE & navigation with rotation actions & its generators need not commute, so it is the one \emph{positional} code that can express a heading-dependent displacement & needing a commuting special case to train at all \\[3pt]
CoPE, CARoPE, GRAPE-AP & \textbf{navigation} & they are monotone, so the companion review predicts they fail on a map task specifically, not generally & matching a signed phase on the torus \\[3pt]
MapFormer / MapEM & transfer across a change of \textbf{structure} --- new topology, new action space, a graph reused across environments & this is TEM's strongest claim and the one $\S$\ref{sec:factorisation} does \emph{not} test: everything measured here is a new instance of the same structure & the structural code failing to survive a topology change, which would confine the factorisation claim to re-instantiation \\[3pt]
any structural code & \textbf{sample efficiency}, not final accuracy & if a where transfers, a new environment should need fewer episodes, not merely reach a higher ceiling. No learning curve has been measured here & equal sample efficiency with a higher asymptote, which would make the separation a capacity effect rather than a transfer one \\
\bottomrule
\end{tabular}
\end{table}

\subsection{What to borrow, specifically, for the what/where separation}
\label{sec:borrow}

If the goal is a cleaner separation of the structural code from the sensory one,
the table says which mechanisms have a part worth taking and which do not.

\textbf{Take CoPE's gate, but not its sign.} MapFormer's separator is a linear
bottleneck: $\Delta$ must reach zero on observations by cancellation in
$W_{\mathrm{out}}W_{\mathrm{in}}$, and nothing makes zero a natural output. CoPE's
is a sigmoid, for which zero is the resting state, and its entire purpose is
deciding \emph{which tokens count}. The evidence that this is the right borrow is
that MapFormer re-derives it unprompted: on contextual counting an unconstrained
increment learns a content gate in $8/8$ seeds, and a magnitude-matched intervention
shows the gate is most of the mechanism rather than a component of it
($\S$\ref{sec:crossover}). But CoPE's gate is $\sigma(\cdot)\in(0,1)$ and therefore
non-negative, which by the companion review builds a clock and not a map. The two
mechanisms answer \emph{different} questions --- a gate says \emph{whether} a token
moves you, a sign says \emph{which way} --- and CoPE has only the first. The
combination to build is a gate on a \emph{signed} increment. One cost is worth
stating: CoPE's gate is computed per query--key \emph{pair}, which is why it has no
$\theta_t$ at all and sits outside the frame; a \emph{per-token} gate is the weaker
but cheap version that keeps the prefix scan.

\textbf{Read PoPE as the what-channel, which reframes MapPoPE.}
$\S$\ref{sec:mappope} treats PoPE's magnitude being \emph{instantaneous} as a
deficiency, because it supplies no second accumulator. Under the factorisation
question it is the opposite: an instantaneous content-driven magnitude carries
sensory information \emph{inside} the positional operator while being structurally
unable to corrupt the structural code, and that is measured --- adding it does not
move the growth exponent ($0.689$ against $0.710$, inside its \textsc{mde}). So
MapPoPE is already a what/where split within the operator: an accumulated signed
phase for the where, a non-accumulated content magnitude for the what. It is also
the best arm measured on the torus task ($0.994$ against $0.967$). The same property
is a defect for a system needing a clock and a virtue for one needing a clean
separation, which is the sort of thing only the frame makes visible.

\textbf{Do not borrow Selective RoPE's gate for this.} It is the obvious candidate
--- a sigmoid on the phase increment --- and it was tested directly. It does not
suppress by token class: on the torus it gives $0.560$ on actions against $0.415$ on
observations, a ratio of $1.35\times$ and nowhere near suppression, and on parity it
is $1.54$--$1.70\times$ in a setting where the gate \emph{hurts}. Whatever it does,
it is not separating what from where.

\textbf{Do not borrow the forget gate for this either.} It sits in the magnitude
slot and is monotone, so it is a clock (the companion review); it suppresses what
is \emph{remembered}, not what updates position. It is the right borrow for a
different problem, listed below.

\textbf{Algebraic PE is the one to take when the structure is known} --- and it is
the most relevant row in the table to a \emph{relational} goal. It maps ``the
algebraic specification of a domain'' to orthogonal operators so that the model
``upholds its desired structural properties'', and it covers sequences, grids,
\textbf{trees, and their compositions}. Two consequences. It expresses structures
MapFormer's abelian $SO(2)$ cannot --- a tree is not a product of circles --- so it
is the natural vehicle for the one prediction in Table~\ref{tab:predictions} that
the factorisation thesis most needs and this work does not test, transfer across a
change of \emph{structure} rather than of instance. And it is the honest oracle for
the separation itself: where MapFormer must \emph{infer} the where from the action
stream, Algebraic PE is handed it, so the gap between them prices what inference
costs.

\textbf{TAPE's idea is already tested here, and it is a null.} TAPE contextualises
the positional encoding ``with sequence content across layers'', refining it at each
update --- i.e. content-driven \emph{re-estimation} of the where along depth. That is
the same object as the refine-$\theta$ variant measured in this project, which
carried and corrected the position estimate at each pass. It does nothing:
$-0.001/-0.011/+0.005$ at training length across three noise levels, with no slope
in noise, and the learned gate settles near zero, so the optimiser \emph{declines}
to refine rather than being unable to. That is a mechanism answer and it applies to
this borrow before it is attempted.

\textbf{HGRN offers a third kind of separation: by timescale, bound to depth.} Its
forget rate carries a lower bound that ``increases monotonically when moving up
layers'', so lower layers model short-range and upper layers long-range structure.
That is neither a slot nor a gate but a third axis on which what and where could be
divided, and it has a concrete form here: MapFormer fixes the same geometric
$\omega$ ladder in every layer, and binding its range to depth instead is a
free change with a published precedent. Untested.

\textbf{Titans gates the other side.} Its memory is written by \emph{surprise} ---
selection on what is worth storing, not on what updates position. A gate on the
where and a gate on the what are complementary, and only the first is discussed
above.

\textbf{And keep the bottleneck.} Since $W_{\mathrm{in}}$ \emph{is} the
identification mechanism, its rank is the channel the whole distinction passes
through, which is why $r=2$'s skewed basis is a fact about the separation and $r=4$
repairs it for $384$ parameters ($\S$\ref{sec:rankmeas}).

\subsection{Combinations worth building}

\begin{table}[htbp]\centering\small
\caption{Combinations the frame identifies, ordered by how much evidence already points at them. None has been built.}
\label{tab:combinations}
\begin{tabular}{@{}p{3.6cm}p{3.9cm}p{3.6cm}p{3.2cm}@{}}
\toprule
combination & what it fixes & the evidence pointing at it & refuted by \\
\midrule
\textbf{signed accumulated phase $+$ accumulated content-dependent magnitude} (MapFormer $+$ a forget gate), in softmax attention & the empty cell: a map \emph{and} a clock at once. Mamba-3 holds it in an SSM; nothing holds it in softmax attention & the forget gate's own accumulator is ballistic ($\alpha = +0.956$) beside the phase's $0.6$, so the clock is already there and unexploited; and it explains why the gate needs a \emph{live} $\lambda$ but not decay & no gain on a task needing both, i.e. recency-dependent navigation, which does not yet exist \\[3pt]
MapPoPE $+$ a forget gate & its actual deficit: it reads as both slots and supplies one answer ($\S$\ref{sec:mappope}) & $\alpha$ unchanged by PoPE ($0.689$ vs $0.710$), so the second accumulator is genuinely absent & $\alpha$ still unmoved, which would mean the magnitude slot cannot host a clock at all \\[3pt]
explicit content gate $+$ signed phase --- as an oracle \textbf{control}, not a proposal & measures what learning the what/where split costs, against a design whose stated purpose is to learn it & on contextual counting an \emph{unconstrained} signed arm learns the gate by itself, and a constant increment on the gated class recovers $0.783$ of $1.000$ ($\S$\ref{sec:crossover}) & the explicit gate not beating the learned one, which would mean the learning is free \\[3pt]
signed phase $+$ non-abelian transport & rotation action spaces, without the allocentric recoding that needs known heading & recoding restores $+0.488$, so the information is sufficient; the question is whether the mechanism can carry it & the parallel-scan loss costing more than the accuracy gained \\[3pt]
rank $r{=}4$ $+$ weight sharing across depth & reliability rather than peak accuracy & measured together: $0.986 \pm 0.020$ with $8/8$ seeds above $0.941$, against $0.416$ for either alone --- the one super-additive pair found here & --- (this one is measured; it is listed because it is the template the others are proposed on) \\
\bottomrule
\end{tabular}
\end{table}

\noindent \textbf{The ordering principle.} Every combination above pairs mechanisms
that occupy \emph{disjoint} coordinates of the companion review. That is the frame's
one practical prescription: mechanisms in the same slot compete and mechanisms in
different slots may compose, and which of the two is happening can be read off the
table before anything is trained. MapPoPE is the cautionary case --- disjoint
coordinates are necessary and not sufficient, because both of its halves answer the
same question.


\section{Conclusion}

Two properties of the content-to-phase map, neither varied deliberately anywhere in
the literature, account for more of the variance on these tasks than the choice of
positional encoding does. The sign costs nothing and decides whether the accumulator
is a map or a clock; the rank costs $384$ parameters and decides whether the map is
well conditioned. Both matter only where position must be reconstructed from the
input, which is why they are invisible on text and why the measurements here are not.

What we could not show is as much of the content. The growth exponent that unifies
the two is nearly collinear with a simpler observable and its route to
length-degradation is refuted rather than established. An explicit separator works
and buys nothing, which corroborates the design it was meant to improve. And the
transfer demonstrated here is to a new \emph{instance} of a structure, not to a new
structure --- which is the claim the cognitive-map literature makes most strongly and
the one still to test.

\end{document}
